
\subsubsection{6.1 Key Findings}

\textbf{FGO's benefit derives from gradient exclusion.} Non-executed tokens receive exactly zero gradient, while executed tokens receive 1.78x stronger signal concentration. This concentrates learning on code that actually contributed to the reward, providing dense supervision without requiring additional reward shaping or curriculum design.

\textbf{Trace collection is reliable but token mapping is imprecise.} The 100\% trace capture rate demonstrates that \texttt{sys.settrace} reliably captures execution information. The 81\% token F1 reflects tokenizer-to-line mapping challenges, not trace collection failure. Improvement paths include AST-based span mapping rather than line-based mapping.

\textbf{Simulation validates mechanism, not production performance.} The 10\% improvement in final pass@1 is observed in simulation-based evaluation. Full training runs are required to validate this finding under realistic conditions.

\subsubsection{6.2 Limitations}

\begin{enumerate}
\item \textbf{Simulation-based efficiency validation}: H-M3 uses simulated learning curves with parametric models, not actual PPO training. Claims about convergence speed or sample efficiency cannot be validated without full training runs.
\end{enumerate}

\begin{enumerate}
\item \textbf{Token mapping precision}: 81\% F1 versus 95\% target due to tokenizer-line misalignment. Some tokens are incorrectly masked or unmasked, introducing noise into the FGO signal.
\end{enumerate}

\begin{enumerate}
\item \textbf{Proof-of-concept scale}: Limited to single-seed evaluation for some experiments. Statistical variance is not fully characterized. Full factorial content-versus-granularity comparison was not executed.
\end{enumerate}

\begin{enumerate}
\item \textbf{Overhead}: Trace collection overhead (21.55x P95) is marginally above the 20x threshold. Production deployment may require optimization such as C-based trace implementations.
\end{enumerate}

\subsubsection{6.3 Unverified Claims}

The following claims from the original hypothesis remain unverified:

\begin{itemize}
\item \textbf{Granularity dominates content}: The $2\times 3$ factorial ANOVA comparing content effects (compile, test, combined) versus granularity effects (standard, FGO) was not executed.
\end{itemize}

\begin{itemize}
\item \textbf{Combined content outperforms single types}: No comparison of feedback content types was conducted.
\end{itemize}

\begin{itemize}
\item \textbf{Faster convergence}: The simulation showed no convergence speedup (steps ratio = 1.0).
\end{itemize}

These claims require full training runs in future work.

\subsubsection{6.4 Scope Conditions}

Results are established for:
\begin{itemize}
\item Python code generation
\item 7B model scale (CodeLlama-7B reference)
\item Function-level tasks (HumanEval, MBPP)
\item PPO algorithm
\end{itemize}

Results may not hold for:
\begin{itemize}
\item Other programming languages (different trace tools required)
\item Repository-level code generation (SWE-bench)
\item Models smaller than 3B or larger than 13B
\item Other RL algorithms (DPO, REINFORCE)
\end{itemize}
